% =====================================================================
%  Capítulo 1 — Introducción
% =====================================================================
\chapter{Introducción}\label{cap:intro}

\entradilla{Antes de empezar, conozca qué es el panel de cotizaciones, qué puede hacer en él y
cómo está organizado este manual.}

\section{¿Qué es el panel de cotizaciones?}\label{sec:que-es}

\textbf{Training~InnovaQuote} es la plataforma web donde los clientes de Training Corporation solicitan
cotizaciones. Está disponible en \ruta{\dominio}. El cliente elige el tipo de servicio (la
\emph{unidad de negocio}), diligencia el \emph{formulario de solicitud} y envía su solicitud.

El \textbf{panel de cotizaciones} es la parte privada de Training~InnovaQuote. Desde allí el equipo de
Training Corporation recibe las solicitudes, envía las propuestas comerciales y mantiene al día
la información que ve el cliente. En el panel usted puede:

\begin{itemize}
  \item Revisar las solicitudes de cotización y enviar la propuesta a cada cliente.
  \item Consultar los clientes que han cotizado y su historial, y exportarlos a Excel.
  \item Administrar las unidades de negocio, las líneas de gestión y los servicios.
  \item Diseñar el formulario de preguntas (árbol de decisión) que responde el cliente al cotizar.
  \item Crear y administrar las cuentas de acceso al panel.
  \item Enviar correos de bienvenida para invitar clientes a la plataforma.
  \item Definir quién recibe las notificaciones y los datos de contacto que llevan los correos.
\end{itemize}

\section{El nombre Training InnovaQuote}\label{sec:nombre}

El nombre de la aplicación une tres palabras, y cada una expresa una parte de lo que es:

\begin{figure}[H]
\centering
\begin{tikzpicture}
  \node[font=\montserrat\fontseries{ub}\selectfont\fontsize{38}{38}\selectfont, inner sep=0pt] (n)
    {\textcolor{noche}{Training}\hspace{0.28em}\textcolor{azul}{Innova}\textcolor{azulvivo}{Quote}};
  \foreach \x/\col/\parte/\icono/\texto in {
      -5.1/noche/Training/F19C/{Respaldo e identidad de Training Corporation},
      0/azul/Innova/F0EB/{Innovación y transformación},
      5.1/azulvivo/Quote/F70E/{Cotizaciones y propuestas comerciales}} {
    \node[tarjeta, draw=\col!60, line width=1.2pt, text width=4.1cm, minimum height=2.9cm,
        align=center, inner ysep=9pt] at (\x,-2.75)
      {{\color{\col}\fontsize{15}{15}\selectfont\ico{\icono}}\\[5pt]
       {\montserrat\fontseries{eb}\selectfont\Large\color{\col}\parte}\\[4pt]
       {\small\color{tinta}\texto}};
  }
\end{tikzpicture}
\caption{Significado del nombre Training InnovaQuote.}\label{fig:nombre}
\end{figure}

\begin{tabla}{colspec={Q[3.2cm,l] X[l]}}
  Palabra & Significado \\
  \textbf{Training} & Respaldo e identidad de Training Corporation. \\
  \textbf{Innova} & Innovación y transformación. \\
  \textbf{Quote} & Cotizaciones y propuestas comerciales. \\
\end{tabla}

\section{¿Quién es el usuario administrador?}\label{sec:quien}

Toda persona que tenga una cuenta en el panel es administrador. El panel no maneja roles ni
permisos distintos: cualquier usuario puede ver y modificar todos los módulos, incluidas las
cuentas de los demás usuarios.

\begin{importante}
Entregue cuentas solo a personas de confianza y cree una cuenta por persona. No comparta su
correo y contraseña con otros: así queda claro quién tiene acceso y puede retirarlo cuando
alguien deje de necesitarlo (ver \seccion{sec:usuarios}).
\end{importante}

\section{Requisitos}\label{sec:requisitos}

\begin{itemize}
  \item Un computador con un navegador web actualizado: Google Chrome, Microsoft Edge,
        Mozilla Firefox o Safari.
  \item Conexión a internet.
  \item La dirección web de Training~InnovaQuote (\ruta{\dominio}) y una cuenta de usuario (correo y
        contraseña) creada por otro administrador.
  \item Acceso al buzón de correo de su cuenta, por si necesita recuperar la contraseña.
\end{itemize}

\begin{consejo}
El panel también funciona en celulares y tabletas (ver \seccion{sec:movil}). Para trabajar con
tablas largas y con el árbol de decisión es más cómodo usar una pantalla amplia.
\end{consejo}

\section{Cómo está organizado este manual}\label{sec:estructura}

\begin{tabla}{colspec={Q[3.2cm,l] X[l]}}
  Parte & Qué encontrará \\
  \Capref{cap:acceso} & Cómo iniciar sesión, recuperar la contraseña y moverse por el panel. \\
  \Capref{cap:operacion} & Las solicitudes de cotización, el envío de propuestas y los clientes. \\
  \Capref{cap:catalogo} & Unidades de negocio, líneas de gestión, servicios y el árbol de decisión. \\
  \Capref{cap:sistema} & Usuarios, correos de bienvenida y configuración. \\
  \Capref{cap:salida} & Cómo ver el sitio público y cerrar sesión. \\
  \Capref{cap:capacidad} & Cuántos usuarios simultáneos soporta Training~InnovaQuote. \\
  \Apref{ap:mensajes} a \Apref{ap:glosario} & Mensajes del sistema, preguntas frecuentes y glosario. \\
\end{tabla}

\section{Cómo leer este manual}\label{sec:convenciones}

Para que sea fácil reconocer en el texto lo que verá en pantalla, el manual usa estas
convenciones:

\begin{tabla}{colspec={Q[5.2cm,l,m] X[l,m]}}
  Así aparece en el manual & Qué significa \\
  \boton{Guardar cambios} & Un botón principal de la pantalla. Haga clic sobre él. \\
  \botonsec{Cancelar} & Un botón secundario, con borde y fondo blanco. \\
  \botonrojo{\icEliminar\ Eliminar} & Un botón que borra información. \\
  \campo{Correo electrónico} & El nombre de un campo, una columna o una etiqueta, tal como aparece en pantalla. \\
  \menu{Catálogo}{Servicios} & Una opción de la barra lateral: primero el grupo, luego la opción. \\
  \msg{Servicio creado correctamente.} & Un mensaje que muestra el sistema, copiado textualmente. \\
  \msg{Pregunta \var{nombre} creada.} & La parte entre ‹ › cambia según el caso; aquí, el texto de la pregunta. \\
  \tecla{Enter} & Una tecla del teclado. \\
  \ruta{/dashboard/servicios} & La dirección de la pantalla, después de \ruta{\dominio}. \\
  \pendiente\ \generada & Insignias de estado, con los mismos colores de la aplicación. \\
\end{tabla}

Además encontrará cuatro tipos de recuadros:

\begin{nota}
Información complementaria que ayuda a entender cómo funciona algo.
\end{nota}

\begin{importante}
Advertencias. Lea con atención: indican acciones que no se pueden deshacer o que afectan a los
clientes.
\end{importante}

\begin{consejo}
Recomendaciones para trabajar más rápido o evitar errores frecuentes.
\end{consejo}

\begin{resultado}
Lo que debe ver en pantalla cuando completa correctamente los pasos.
\end{resultado}

\section{Mapa del panel}\label{sec:mapa}

Todas las pantallas del panel se abren desde la barra lateral azul, organizada en tres grupos.
Al final de la barra están las opciones para ver el sitio público y cerrar sesión. La
\figref{fig:mapa-panel} resume la estructura y la página de este manual donde se explica cada
módulo.

\newcommand{\modmapa}[3]{{\color{azul}\ico{#1}}\hspace{0.5em}{\fontseries{sb}\selectfont\footnotesize #2}%
  \hfill{\scriptsize\color{pizarra}\pageref*{#3}}}

\begin{figure}[H]
\centering
\begin{tikzpicture}[
    grupo/.style={fill=noche, text=white, font=\montserrat\bfseries\small,
      minimum width=4.7cm, minimum height=0.85cm, inner sep=0pt},
    modulo/.style={draw=linea, line width=0.9pt, fill=white, text=tinta, font=\small,
      minimum height=0.9cm, inner xsep=8pt, text width=4.2cm, align=left, rounded corners=3pt},
    conector/.style={draw=azulvivo, line width=0.9pt},
  ]
  \node[fill=azul, text=white, font=\montserrat\bfseries, minimum width=15.6cm, minimum height=1.05cm]
    (panel) at (0,0) {\ico{F015}\hspace{0.6em}Panel de inicio\hspace{0.6em}{\fontseries{m}\selectfont\small\color{cielo}(pág.~\pageref*{sec:panel-inicio})}};
  \foreach \x/\nombre/\id in {-5.45/Operación/op, 0/Catálogo/ca, 5.45/Sistema/si} {
    \node[grupo] (\id) at (\x,-1.5) {\nombre};
    \draw[conector] (\x,-0.525) -- (\id.north);
  }
  \node[modulo, below=0.3cm of op] (op1) {\modmapa{F70E}{Cotizaciones}{sec:cotizaciones}};
  \node[modulo, below=0.2cm of op1] (op2) {\modmapa{F2B9}{Clientes}{sec:clientes}};
  \node[modulo, below=0.3cm of ca] (ca1) {\modmapa{F0B1}{Unidades de negocio}{sec:unidades}};
  \node[modulo, below=0.2cm of ca1] (ca2) {\modmapa{F126}{Líneas de gestión}{sec:lineas}};
  \node[modulo, below=0.2cm of ca2] (ca3) {\modmapa{F0AD}{Servicios}{sec:servicios}};
  \node[modulo, below=0.2cm of ca3] (ca4) {\modmapa{F0E8}{Árbol de decisión}{sec:arbol}};
  \node[modulo, below=0.3cm of si] (si1) {\modmapa{F0C0}{Usuarios}{sec:usuarios}};
  \node[modulo, below=0.2cm of si1] (si2) {\modmapa{F0E0}{Correos de bienvenida}{sec:bienvenida}};
  \node[modulo, below=0.2cm of si2] (si3) {\modmapa{F013}{Configuración}{sec:configuracion}};
  % Pie de la barra lateral
  \draw[linea, line width=0.9pt, dash pattern=on 4pt off 3pt] (-7.8,-7.15) -- (7.8,-7.15);
  \node[anchor=south west, inner sep=0pt, font=\montserrat\bfseries\scriptsize, text=pizarra]
    at (-7.8,-7.0) {\addfontfeatures{LetterSpace=12}AL FINAL DE LA BARRA LATERAL};
  \node[modulo, text width=7.1cm, anchor=north west] (pub) at (-7.8,-7.45)
    {\modmapa{F08E}{Ver sitio público}{sec:sitio-publico}};
  \node[modulo, text width=7.1cm, anchor=north east] (out) at (7.8,-7.45)
    {\modmapa{F2F5}{Cerrar sesión}{sec:cerrar-sesion}};
\end{tikzpicture}
\caption{Estructura del panel de cotizaciones.}\label{fig:mapa-panel}
\end{figure}

\section{Recorrido rápido}\label{sec:recorrido}

Este es el recorrido completo de una jornada de trabajo en el panel, desde que entra hasta que
sale. Cada paso remite a la parte del manual donde se explica en detalle.

\begin{figure}[H]
\centering
\begin{tikzpicture}[
    etapa/.style={signal, signal to=east, signal from=west, signal pointer angle=125,
      minimum height=1.9cm, minimum width=3.2cm, inner xsep=4pt, text width=2.1cm,
      align=center, font=\montserrat\bfseries\scriptsize},
  ]
  \foreach \i/\fondo/\letra/\ic/\tx/\etiq in {
      0/noche/white/F084/Iniciar sesión/sec:iniciar-sesion,
      1/azulhondo/white/F015/Panel de inicio/sec:panel-inicio,
      2/azul/white/F70E/Trabajar en los módulos/cap:operacion,
      3/azulvivo/white/F08E/Ver sitio público/sec:sitio-publico,
      4/cielo/noche/F2F5/Cerrar sesión/sec:cerrar-sesion} {
    \node[etapa, fill=\fondo, text=\letra] at (\i*3.12,0)
      {{\fontsize{12}{12}\selectfont\ico{\ic}}\\[3pt]\tx};
    \node[font=\scriptsize, text=pizarra] at (\i*3.12,-1.3) {pág.~\pageref*{\etiq}};
  }
\end{tikzpicture}
\caption{Recorrido de una sesión de trabajo en el panel.}\label{fig:recorrido}
\end{figure}

\enlargethispage{4\baselineskip}
\begin{pasos}
  \item \textbf{Inicie sesión} con su correo y contraseña (\seccion{sec:iniciar-sesion}).
  \item En el \textbf{Panel de inicio}, elija el módulo con el que va a trabajar
        (\seccion{sec:panel-inicio}).
  \item Abra \menu{Operación}{Cotizaciones} y revise las solicitudes en estado \pendiente\
        (\seccion{sec:cotizaciones}).
  \item Abra cada solicitud, lea las respuestas del cliente y descargue los documentos que adjuntó.
  \item Cuando la propuesta esté lista, envíela al cliente con \boton{\icCargar\ Enviar propuesta}
        (\seccion{sec:enviar-propuesta}).
  \item Consulte \menu{Operación}{Clientes} cuando necesite el historial de un cliente o un
        reporte en Excel (\seccion{sec:clientes}).
  \item Si cambia la oferta de servicios, actualice el catálogo y el árbol de decisión
        (\capref{cap:catalogo}).
  \item Revise cómo lo ve el cliente con \textbf{Ver sitio público} y, al terminar,
        \textbf{cierre sesión} (\capref{cap:salida}).
\end{pasos}
